\section{ばね・振り子}

\subsection{平成27年度 問1}
\label{q:h27-reg-applied-01}

\noindent\textbf{原典：} sources/H27/17\_applied\_science.pdf、PDF物理1頁。

\textbf{公式正答：} 4（single）。

\paragraph{問題の要旨} ばね定数1.0 N/mの軽いばねに質量8.0 kgをつるした鉛直単振動の周期を求める。

\paragraph{問われる核心} 復元力の線形化と固有角振動数

\paragraph{必要知識} \(m\ddot x+kx=0\)なら\(T=2\pi\sqrt{m/k}\)。単振り子は小角近似で\(\ddot\theta+(g/\ell)\theta=0\)。

\paragraph{考え方（査読待ち初稿）} \(T=2\pi\sqrt{m/k}=2\pi\sqrt8\simeq17.8\) sで④。重力は平衡位置をずらすが、平衡位置からの振動周期には入らない。

\paragraph{正答を選ぶ理由} 独立に得た条件適合値・組合せは選択肢4である。導出した周期に最も近い。角周波数と周期を区別して2πを含める。

\paragraph{選択肢の検討（査読待ち初稿）}\mbox{}\par

\noindent\textbf{1：不採用} 1.4 s。指定条件を振動周期式へ代入した値と一致しない。\par

\noindent\textbf{2：不採用} 2.8 s。指定条件を振動周期式へ代入した値と一致しない。\par

\noindent\textbf{3：不採用} 8.8 s。指定条件を振動周期式へ代入した値と一致しない。\par

\noindent\textbf{4：採用} 18 s。導出した周期に最も近い。角周波数と周期を区別して2πを含める。\par

\noindent\textbf{5：不採用} 26 s。指定条件を振動周期式へ代入した値と一致しない。\par

\paragraph{再利用する考え方} 復元力の線形化と固有角振動数を、条件と式の対応から説明できること。

\paragraph{共通理論} \hyperref[theory:mechanics:oscillation]{第\ref*{theory:mechanics:oscillation}節}

\par 難易度 1（校準初稿）、優先度未評価。

\subsection{平成29年度 問1}
\label{q:h29-reg-applied-01}

\noindent\textbf{原典：} sources/H29/17\_applied\_science.pdf、PDF物理1頁。

\textbf{公式正答：} 3（single）。

\paragraph{問題の要旨} 長さ0.1 mのひもに質量0.1 kgの小さいおもりをつけた小角単振り子の周期を求める。

\paragraph{問われる核心} 復元力の線形化と固有角振動数

\paragraph{必要知識} \(m\ddot x+kx=0\)なら\(T=2\pi\sqrt{m/k}\)。単振り子は小角近似で\(\ddot\theta+(g/\ell)\theta=0\)。

\paragraph{考え方（査読待ち初稿）} 小角近似\(\sin\theta\simeq\theta\)から\(T=2\pi\sqrt{\ell/g}\)。\(T=2\pi\sqrt{0.1/9.8}\simeq0.635\) sで③。質量は消える。

\paragraph{正答を選ぶ理由} 独立に得た条件適合値・組合せは選択肢3である。導出した周期に最も近い。角周波数と周期を区別して2πを含める。

\paragraph{選択肢の検討（査読待ち初稿）}\mbox{}\par

\noindent\textbf{1：不採用} 0.1 s。指定条件を振動周期式へ代入した値と一致しない。\par

\noindent\textbf{2：不採用} 0.3 s。指定条件を振動周期式へ代入した値と一致しない。\par

\noindent\textbf{3：採用} 0.6 s。導出した周期に最も近い。角周波数と周期を区別して2πを含める。\par

\noindent\textbf{4：不採用} 1.0 s。指定条件を振動周期式へ代入した値と一致しない。\par

\noindent\textbf{5：不採用} 6.0 s。指定条件を振動周期式へ代入した値と一致しない。\par

\paragraph{再利用する考え方} 復元力の線形化と固有角振動数を、条件と式の対応から説明できること。

\paragraph{共通理論} \hyperref[theory:mechanics:oscillation]{第\ref*{theory:mechanics:oscillation}節}

\par 難易度 1（校準初稿）、優先度未評価。

\subsection{令和7年度 問2}
\label{q:r07-reg-applied-02}

\noindent\textbf{原典：} sources/R07/17\_applied\_science.pdf、PDF物理1頁。

\textbf{公式正答：} 1（single）。

\paragraph{問題の要旨} 質量0.80 kgを鉛直ばねにつるすと4.9 cm伸びた。平衡位置から引いて放す単振動の周期を求める。

\paragraph{問われる核心} 復元力の線形化と固有角振動数

\paragraph{必要知識} \(m\ddot x+kx=0\)なら\(T=2\pi\sqrt{m/k}\)。単振り子は小角近似で\(\ddot\theta+(g/\ell)\theta=0\)。

\paragraph{考え方（査読待ち初稿）} 静的釣合い\(k\Delta\ell=mg\)から\(k=0.80\times9.8/0.049=160\) N/m。\(T=2\pi\sqrt{m/k}=2\pi\sqrt{0.80/160}\simeq0.444\) sで①。cmをmへ変換してから計算する。

\paragraph{正答を選ぶ理由} 独立に得た条件適合値・組合せは選択肢1である。導出した周期に最も近い。角周波数と周期を区別して2πを含める。

\paragraph{選択肢の検討（査読待ち初稿）}\mbox{}\par

\noindent\textbf{1：採用} \(4.4\times10^{-1}\) s。導出した周期に最も近い。角周波数と周期を区別して2πを含める。\par

\noindent\textbf{2：不採用} \(8.9\times10^{-1}\) s。指定条件を振動周期式へ代入した値と一致しない。\par

\noindent\textbf{3：不採用} 1.0 s。指定条件を振動周期式へ代入した値と一致しない。\par

\noindent\textbf{4：不採用} 1.8 s。指定条件を振動周期式へ代入した値と一致しない。\par

\noindent\textbf{5：不採用} 2.0 s。指定条件を振動周期式へ代入した値と一致しない。\par

\paragraph{再利用する考え方} 復元力の線形化と固有角振動数を、条件と式の対応から説明できること。

\paragraph{共通理論} \hyperref[theory:mechanics:oscillation]{第\ref*{theory:mechanics:oscillation}節}

\par 難易度 2（校準初稿）、優先度未評価。

\section{波動・音}

\subsection{平成23年度 問3}
\label{q:h23-reg-applied-03}

\noindent\textbf{原典：} sources/H23/17\_applied\_science.pdf、PDF物理1頁。

\textbf{公式正答：} 2（single）。

\par\textit{原典本文の校正と独自解説は未作成。分類のみの登録項目。}

\subsection{平成24年度 問2}
\label{q:h24-reg-applied-02}

\noindent\textbf{原典：} sources/H24/17\_applied\_science.pdf、PDF物理1頁。

\textbf{公式正答：} 4（single）。

\par\textit{原典本文の校正と独自解説は未作成。分類のみの登録項目。}

\subsection{平成25年度 問2}
\label{q:h25-reg-applied-02}

\noindent\textbf{原典：} sources/H25/17\_applied\_science.pdf、PDF物理1頁。

\textbf{公式正答：} 3（single）。

\par\textit{原典本文の校正と独自解説は未作成。分類のみの登録項目。}

\subsection{平成25年度 問5}
\label{q:h25-reg-applied-05}

\noindent\textbf{原典：} sources/H25/17\_applied\_science.pdf、PDF物理3頁。

\textbf{公式正答：} 2（single）。

\par\textit{原典本文の校正と独自解説は未作成。分類のみの登録項目。}

\subsection{平成29年度 問4}
\label{q:h29-reg-applied-04}

\noindent\textbf{原典：} sources/H29/17\_applied\_science.pdf、PDF物理2頁。

\textbf{公式正答：} 5（single）。

\par\textit{原典本文の校正と独自解説は未作成。分類のみの登録項目。}

\subsection{平成30年度 問5}
\label{q:h30-reg-applied-05}

\noindent\textbf{原典：} sources/H30/17\_applied\_science.pdf、PDF物理2頁。

\textbf{公式正答：} 2（single）。

\par\textit{原典本文の校正と独自解説は未作成。分類のみの登録項目。}

\subsection{令和元年度 問5}
\label{q:r01-reg-applied-05}

\noindent\textbf{原典：} sources/R01/17\_applied\_science.pdf、PDF物理2頁。

\textbf{公式正答：} 2（single）。

\par\textit{原典本文の校正と独自解説は未作成。分類のみの登録項目。}

\subsection{令和4年度 問5}
\label{q:r04-reg-applied-05}

\noindent\textbf{原典：} sources/R04/17\_applied\_science.pdf、PDF物理2頁。

\textbf{公式正答：} 2（single）。

\par\textit{原典本文の校正と独自解説は未作成。分類のみの登録項目。}

\subsection{令和5年度 問5}
\label{q:r05-reg-applied-05}

\noindent\textbf{原典：} sources/R05/17\_applied\_science.pdf、PDF物理2頁。

\textbf{公式正答：} 3（single）。

\par\textit{原典本文の校正と独自解説は未作成。分類のみの登録項目。}

\subsection{令和7年度 問5}
\label{q:r07-reg-applied-05}

\noindent\textbf{原典：} sources/R07/17\_applied\_science.pdf、PDF物理2頁。

\textbf{公式正答：} 4（single）。

\par\textit{原典本文の校正と独自解説は未作成。分類のみの登録項目。}
